
\section{Conclusion}

We introduced SECDA-LLM, a novel platform that simplifies the creation of specialized hardware accelerators for LLM inference on resource-constrained edge devices. 
% SECDA-LLM integrates the SECDA methodology within \LPP and provides tools to streamline accelerator design development.
% By integrating SECDA environment with \LPP, developers can effectively co-design new FPGA-based accelerators for LLMs with ease using the SECDA design flow.
SECDA-LLM integrates the SECDA methodology within \LPP, enabling developers to access the SECDA tools (e.g., AXI-API, profiler), which can be used to effectively co-design new FPGA-based accelerators for LLMs with ease.
As a case study, we presented a quantized MatMul accelerator design that optimizes LLM inference (by $11\times$ over the dual-core Arm NEON-based CPU execution) for the TinyLlama model.
Future work will expand SECDA-LLM into an open-source platform for collaborative development and continuous improvement of LLMs' performance on resource-constrained edge devices.




% We introduced SECDA-LLM, a novel platform that simplifies the creation of new FPGA-based hardware accelerators for edge LLM inference using the hardware/software co-design SECDA methodology within \LPP.
% SECDA-LLM removes the initial setup within the \LPP framework and tools to streamline accelerator design development.
% By integrating the SECDA tools with \LPP, developers can effectively co-design new accelerators for LLMs with ease using the SECDA design flows.

% Our approach using the SECDA methodology reduces development costs by minimizing synthesis iterations and replacing them with simulation iterations using SystemC for fine-grained benchmarking and co-verification.
% Furthermore, designs can be directly deployed onto FPGAs from simulations, eliminating the need for re-implementation.
% As a case study, we presented a quantized MatMul-based accelerator design that optimizes LLM inference for the TinyLlama model.
% Future work will expand SECDA-LLM into an open-source platform for collaborative development and continuous improvement of LLMs' performance on edge devices.





% Our approach using the SECDA methodology reduces development costs by minimizing synthesis iterations and replacing them with simulation iterations using SystemC for fine-grained benchmarking and co-verification.
% Furthermore, designs can be directly deployed onto FPGAs from simulations, eliminating the need for re-implementation.
% As a case study, we presented a quantized MatMul-based accelerator design that optimizes LLM inference for the TinyLlama model.
% Future work will expand SECDA-LLM into an open-source platform for collaborative development and continuous improvement of LLMs' performance on edge devices.




% Long conclusion
% In this paper, we presented SECDA-LLM, a new platform that increases the ease of developing new hardware accelerators for edge LLM inference, utilizing the hardware/software co-design SECDA methodology. 
% SECDA-LLM provides the initial environment setup within the \LPP inference framework as well as a set of tools and utilities to aid in the development of accelerator designs, providing the infrastructure to initialize and follow the SECDA environment within \LPP
% As a result, developers are able to co-design new accelerator designs for LLMs models, bypassing many of the initial setup costs such as valuable developer time.


% We reduce development costs by minimizing the number of synthesis iterations, which are replaced by simulation iterations.
% Simulating the accelerator behavior using SystemC, allows fine-grained benchmarking and co-verification of accelerator components. 
% Furthermore, we can directly map our simulated designs onto an FPGA without re-implementation.
% As a case study, we proposed quantised MatMul-based accelerator design for optimizing the inference of LLM models.
% We evaluate our accelerator on the TinyLlama model, where we outperform the CPU baseline by up to \uptooverallX.
% As future work, we plan to develop the SECDA-LLM platform into an open-source environment for accelerator developers collaborate and contribute to improving performance of LLMs for edge devices.